\documentclass[11pt]{article}
\usepackage[margin=1in]{geometry}
\usepackage{amsmath,amssymb,amsthm,mathtools}
\usepackage[T1]{fontenc}
\usepackage{lmodern}
\usepackage{microtype}
\usepackage[hidelinks,breaklinks]{hyperref}
\usepackage{enumitem}
\setlist{itemsep=3pt,topsep=5pt}
\newtheorem{theorem}{Theorem}[section]
\newtheorem{proposition}[theorem]{Proposition}
\newtheorem{lemma}[theorem]{Lemma}
\newtheorem{corollary}[theorem]{Corollary}
\theoremstyle{definition}
\newtheorem{assumption}[theorem]{Imported input}
\newtheorem{definition}[theorem]{Definition}
\theoremstyle{remark}
\newtheorem{remark}[theorem]{Remark}
\newcommand{\dd}{\,\mathrm d}
\newcommand{\Rea}{\operatorname{Re}}
\newcommand{\Ima}{\operatorname{Im}}
\DeclareMathOperator{\Log}{Log}
\newcommand{\Hspec}{\mathsf H}
\newcommand{\Abulk}{A_{\mathrm{bulk}}}
\newcommand{\one}{\mathbf1}
\newcommand{\Sch}{\mathcal S}
\title{Enlarged-state transport and visibility\texorpdfstring{\\}{ }for the Newman heat flow}
\author{Drafted for Edward Baker}
\date{10 October 2026}

\begin{document}
\maketitle
\begin{abstract}
We organize an enlarged-state approach to real multiple zeros of the
Riemann heat flow. A complete finite prime state obeys both commuting
spectral dynamics and noncommuting divisor transport. Anchoring its
arithmetic source produces a strictly positive operator with an exact
source-response formula. A conditional adjoint-transport certificate then
bounds the joint value and slope using only the imported first-jet
approximation, and therefore applies to every multiplicity on its domain.
We prove an eventual exclusion criterion conditional on uniform residual
and multiplier bounds; on the stated shrinking sector a sufficient
multiplier growth exponent is less than $5/8$.
A continuous theta wavefunction gives a chord equation and a coupled
signed gradient overlap. A retained angular theta field supplies an exact
differential-shift identity with an endpoint source. Logarithmic reflection
organizes the natural cutoff as matched interior and exterior channels.
These are exact reductions and conditional criteria. No useful uniform
multiplier construction, theta-specific overlap sign, new collision
exclusion, or proof of the Riemann hypothesis is established.
\end{abstract}

\section{The genuine flow and the paid observation}
\label{sec:flow}

The normalization throughout is
\begin{align}
 H_t(z)&=\int_0^\infty e^{tu^2}\Phi(u)\cos(zu)\dd u,
 &H_0(z)&=\tfrac18\xi(\tfrac12+iz/2),\label{eq:heat}\\
 \Phi(u)&=\sum_{n\ge1}
 (2\pi^2n^4e^{9u}-3\pi n^2e^{5u})e^{-\pi n^2e^{4u}}
 \quad(u\ge0).\label{eq:theta-kernel}
\end{align}
Here $\xi(s)=s(s-1)\pi^{-s/2}\Gamma(s/2)\zeta(s)/2$ is the completed
Riemann function. Let $\Phi_e$ denote the smooth even extension of
$\Phi$. Theta transformation gives this evenness. On $u\ge0$ every
summand in \eqref{eq:theta-kernel} is positive, and the decay is faster
than every fixed Gaussian and exponential weight. These standard
normalizations and the effective heat-flow framework are recorded in
\cite{polymath}; the classical real-zero threshold originates in
\cite{debruijn,newman}.

\begin{proposition}[Physical equation and observation domain]
\label{prop:physical}
For real $t$ the function $H_t$ is entire in $z$, is real on the real
axis, and satisfies
\begin{equation}
 \partial_tH_t=-\partial_z^2H_t.
 \label{eq:newman-sign}
\end{equation}
Its derivatives of every fixed order may be taken under the integral,
locally uniformly in bounded complex $z$ and bounded real $t$.
\end{proposition}
\begin{proof}
The theta envelope dominates $e^{tu^2+|\Ima z|u}$ times each fixed
polynomial in $u$. Dominated differentiation is therefore available on
compact parameter sets. Time differentiation inserts $u^2$, while
$\partial_z^2\cos(zu)=-u^2\cos(zu)$. Reflection gives the stated real
symmetry.
\end{proof}

The local target is $H_t(x)=H_t'(x)=0$ at real $x$. Every prime on a
spatial function means differentiation at fixed time. In the finite
model the integer cutoff is also fixed; derivatives never follow a
shrinking-sector scaling curve. A direct exclusion of this pair excludes
all real multiplicities at least two. Relating uniform positive-time
collision exclusion to RH additionally uses the genuine threshold,
localization, and splitting inputs of the parent investigation
\cite{heatnotes}; no such global transfer is newly proved here.

For motivation we import the parent's finite-threshold collision reduction.
If the genuine de Bruijn--Newman threshold $\Lambda$ were positive,
the bound $\Lambda\le0.22$ and the effective large-height reality and
simplicity statements of \cite[Theorem~1.5(i),(ii)]{polymath} apply
uniformly on a time neighborhood bounded away from zero. Thus zeros
outside a fixed compact height region cannot cause loss of all-realness
there. Continuation of simple real zeros
\cite[Proposition~3.1(i)]{polymath} and compact zero control in the
remaining region force a finite real multiple zero at the threshold.
This uses the imported global heat-flow reduction and its compactness
hypotheses, rather than the local identities proved below. The project
seeks a visibility mechanism on restricted domains first; the full
threshold-covering theorem remains open.

\subsection{The full fixed-cutoff dictionary}

For a complex spatial argument $z$ set $s=(1-iz)/2$. On the large-height
neighborhoods avoiding the principal logarithm cuts, define
\begin{align}
 m_0(s)&=\Log s+\Log(s-1)-\tfrac{s}{2}\log\pi
  +\log\frac{\sqrt{2\pi}}{16}
  +(\tfrac{s}{2}-\tfrac12)\Log\tfrac{s}{2}-\tfrac{s}{2},\label{eq:m0}\\
 \alpha(s)&=m_0'(s)=\frac1{2s}+\frac1{s-1}
                  +\frac12\Log\frac{s}{2\pi},
 &m_t(s)&=m_0(s)+\frac t4\alpha(s)^2.\nonumber
\end{align}
With a fixed positive integer $N$ put
\begin{align}
 P_{t,N}(s)&=\sum_{n\le N}
  \exp\{t\log^2n/4-(s+t\alpha(s)/2)\log n\},\nonumber\\
 G_{t,N}(z)&=e^{m_t(s)}P_{t,N}(s)
            +e^{m_t(1-s)}P_{t,N}(1-s),\nonumber\\
 A_t(z)&=\exp\{(m_t(s)+m_t(1-s))/2\},
 &\theta(t,z)&=\frac{m_t(s)-m_t(1-s)}{2i},\label{eq:normalizer}\\
 Q_t(z)&=H_t(z)/A_t(z),
 &F_{t,N}(z)&=G_{t,N}(z)/A_t(z).\nonumber
\end{align}
The normalizer is nonvanishing; on the real axis it is positive and
$Q_t,F_{t,N},\theta$ are real. Consequently
\begin{equation}
 H_t(x)=H_t'(x)=0\quad\Longleftrightarrow\quad
 Q_t(x)=Q_t'(x)=0.
 \label{eq:normalizer-zero}
\end{equation}

\begin{assumption}[Complete holomorphic approximation]
\label{input:disk}
At a chosen real center in the sector
\begin{equation}
 0<t\le1/20,\quad 1\le\kappa\le3/2,\quad
 L=\kappa/t,\quad x=4\pi e^L,\quad
 N=\left\lfloor\sqrt{e^L+t/16}\right\rfloor,
 \label{eq:sector}
\end{equation}
freeze $N$ on the complex disk and choose a strictly positive $\eta_N$
in the complete bound
\begin{equation}
 |Q_t(z)-F_{t,N}(z)|\le\eta_N\quad(|z-x|\le L^{-1}),
 \qquad \eta_N\le5e^{-\kappa(\kappa+4)/(16t)}.
 \label{eq:disk}
\end{equation}
This includes normalization, reflection, and natural-cutoff changes in
the disk. It is the precise interface imported from the parent heat
investigation using the approximation machinery of \cite{polymath};
it is not asserted as a new theorem of this paper, or as an unqualified
verbatim theorem of that source. No smaller block inherits this input.
\end{assumption}

Cauchy's formula gives
$|Q_t^{(j)}(x)-F_{t,N}^{(j)}(x)|\le j!L^j\eta_N$.
We use the full disk to retain the correlation between its first two jets.

\begin{lemma}[Correlated candidate and its support]
\label{lem:paid-body}
Under Input~\ref{input:disk}, any genuine joint zero satisfies
\begin{equation}
 \left(\frac{F_{t,N}}{\eta_N}\right)^2+
 \left|\frac{F_{t,N}'}{L\eta_N}\right|\le1.
 \label{eq:body}
\end{equation}
For $D,B\ge0$, define the continuous support function
\begin{equation}
 h(D,B)=
 \begin{cases}
 D+B^2/(4D),&D>0,\ B\le2D,\\
 B,&D=0\text{ or }B\ge2D.
 \end{cases}
 \label{eq:support}
\end{equation}
For $\mathcal C_0=\{(s,v):s^2+|v|\le1\}$,
$\sup_{\mathcal C_0}|Bs+Dv|=h(D,B)$.
\end{lemma}
\begin{proof}
Apply Schwarz--Pick to
$g(\omega)=(Q_t-F_{t,N})(x+\omega/L)/\eta_N$ on the unit disk.
At a joint zero its real first jet is
$g(0)=-F_{t,N}/\eta_N$,
$g'(0)=-F_{t,N}'/(L\eta_N)$, and
$|g'(0)|\le1-|g(0)|^2$. This proves \eqref{eq:body}.
For the support, align signs and maximize
$B y+D(1-y^2)$ over $0\le y\le1$.
For $D>0$ the maximizer is $\min\{B/(2D),1\}$; the degenerate case is
direct.
\end{proof}

Writing $S=\sum q_n=u+iv$, so that $F=2u$, the complete current
$\mathcal J=-\Ima(S'\overline S)=u'v-uv'$ obeys at a candidate
\begin{equation}
 |\mathcal J|\le\frac{\eta_N}{2}
                   h(L|v|,|v'|).
 \label{eq:current-support}
\end{equation}
This is a support consequence of \eqref{eq:body}; useful current
exclusion still requires a strict reverse estimate from the full state.

\section{Finite spectral dynamics and arithmetic transport}
\label{sec:finite}

At a real center write $\alpha=\alpha_r+i\alpha_i$,
$\alpha'=U+iV$, with the latter derivative taken in $s$. Define
\begin{align}
 \sigma&=\tfrac12+t\alpha_r/2,&T&=(x-t\alpha_i)/2,
 &\zeta&=\sigma-iT,\nonumber\\
 c&=\tfrac12(1+tU/2),&d&=tV/4,
 &\Omega&=\tfrac12\Rea\{\alpha(1+t\alpha'/2)\}.\label{eq:physical-data}
\end{align}
Then $\theta_x=-\Omega$, $\zeta_x=d-ic$, and the complete physical
coefficients and reduction are
\begin{equation}
 q_n=e^{t\ell_n^2/4-\sigma\ell_n+i(\theta+T\ell_n)},
 \quad\ell_n=\log n,\quad
 w_n=|q_n|,\quad F=2\Rea\sum_{n\le N}q_n.
 \label{eq:q}
\end{equation}
The physical derivative includes drift:
$q_n'=(-i\Omega+(-d+ic)\ell_n)q_n$.

On $\mathbb C^N$ define $\Hspec e_n=\ell_ne_n$ and
$\one_N=\sum_{n\le N}e_n$. The spectral state is exactly
\begin{equation}
 q=e^{i\theta}e^{t\Hspec^2/4-\zeta\Hspec}\one_N.
 \label{eq:state}
\end{equation}
For independent variables $\tau,\zeta$, the carrier-free vector
$\psi_N=e^{\tau\Hspec^2/4-\zeta\Hspec}\one_N$ satisfies
\begin{equation}
 \partial_\tau\psi_N=\Hspec^2\psi_N/4
                     =\partial_\zeta^2\psi_N/4.
 \label{eq:spectral-heat}
\end{equation}
Equivalently the prime polynomial
$P_\tau(W)=\sum_{n\le N}e^{\tau\ell_n^2/4}W^{\nu(n)}$ obeys
$\partial_\tau P_\tau=\mathcal D^2P_\tau/4$, where
$\mathcal D=\sum_{p\le N}(\log p)W_p\partial_{W_p}$.
Unique prime factorization gives $\mathcal DW^{\nu(n)}=\ell_nW^{\nu(n)}$.

The physical pullback has connection matrices
\begin{equation}
 q_x=(i\theta_xI-\zeta_x\Hspec)q,
 \qquad q_t=(i\theta_tI+\Hspec^2/4-\zeta_t\Hspec)q.
 \label{eq:pullback}
\end{equation}
They commute, and their mixed derivatives agree, so the connection is
flat. Every constant diagonal phase matrix commutes with these equations.
In particular $q_n^\chi=\chi(n)q_n$ for a constant completely multiplicative unit
twist obeys the same local diagonal dynamics. Flatness therefore does
not select the physical initial preparation. The finite scalar $F$ is
not an autonomous Newman solution: its physical carrier, spectral
coordinates, and normalization move, and the exact finite heat residual
in program 09 remains present \cite{torusnotes}.

For $t>0$ an auxiliary realization is
\begin{equation}
 e^{t\Hspec^2/4-\zeta\Hspec}\one_N
 =\frac1{\sqrt{\pi t}}\int_{\mathbb R}e^{-\lambda^2/t}
             e^{-(\zeta-\lambda)\Hspec}\one_N\dd\lambda.
 \label{eq:gaussian}
\end{equation}
Completion of the square proves it. The same real $\lambda$ couples all
prime channels. Separate slices are not candidates; positivity of this
measure does not give a joint observation lower bound. All expressions
above are finite. Positive-time exponentiation is unbounded on an
unrestricted infinite spectral span, and $\sum_n e^{t\log^2n/4}n^{-s}$
diverges for each fixed $s$ when $t>0$.

\subsection{Noncommuting divisor transport}

For a prime $p\le N$ let $S_pe_n=e_{pn}$ if $pn\le N$, and zero
otherwise. With $a_p=\log p$,
\begin{equation}
 [\Hspec,S_p]=a_pS_p,\qquad
 [\Hspec^2,S_p]=2a_pS_p\Hspec+a_p^2S_p.
 \label{eq:prime-commutators}
\end{equation}
Define $B_p:\mathbb C^N\to\mathbb C^{\lfloor N/p\rfloor}$ by
\begin{equation}
 (B_pv)_n=v_{pn}-r_{p,n}v_n,
 \quad r_{p,n}=e^{t(2a_p\ell_n+a_p^2)/4-\zeta a_p}.
 \label{eq:edges}
\end{equation}

\begin{proposition}[Fixed physical transport and its tangent]
\label{prop:transport}
The prescribed state satisfies $B_pq=0$ and $q_1=e^{i\theta}$.
These equations uniquely specify it on the complete cutoff.
Moreover
\begin{align}
 (B_pq_x)_n&=-\zeta_xa_pq_{pn},\nonumber\\
 (B_pq_t)_n&=(a_p\ell_n/2+a_p^2/4-\zeta_ta_p)q_{pn}.\label{eq:tangent}
\end{align}
A completely multiplicative unit twist has defect
\begin{equation}
 (B_pq^\chi)_n=\chi(n)(\chi(p)-1)q_{pn}.
 \label{eq:twist-defect}
\end{equation}
\end{proposition}
\begin{proof}
The ratio $q_{pn}/q_n$ from \eqref{eq:q} is exactly $r_{p,n}$.
Every $n\le N$ reaches $1$ by prime division within the cutoff, proving
uniqueness. Differentiating $B_pq=0$ gives
$B_pq_x=(r_{p,n})_xq_n$, with
$(\log r_{p,n})_x=-\zeta_xa_p$; its time derivative gives the second
formula. The carrier cancels in these differences. Complete
multiplicativity gives \eqref{eq:twist-defect}.
\end{proof}

Fixed edge values and the known source retain information lost by the
diagonal connection. Abstract flatness around prime cycles is insufficient:
$r_{p,qn}r_{q,n}=r_{q,pn}r_{p,n}$ still holds if all edges are altered
by a completely multiplicative unit twist. An exact dilation/source-shift identity is
\begin{equation}
 \sum_{n\le N/p}q_{pn}
 =e^{-\zeta a_p+ta_p^2/4}e^{i\theta}
   \sum_{n\le N/p}e^{t\ell_n^2/4-(\zeta-ta_p/2)\ell_n}.
 \label{eq:shifted-source}
\end{equation}
The shifted state has a smaller literal cutoff and receives no separate
candidate equation.

\section{Anchored source response and visibility}
\label{sec:visibility}

Stack all edge operators into $B$. This includes every prime and every
edge retained in \eqref{eq:edges}. The graph is connected, so
$\ker B=\mathbb Cq$.

\begin{proposition}[Positive source-response operator]
\label{prop:source}
With $e_1$ the $n=1$ coordinate source,
\begin{equation}
 \Abulk=B^*B+e_1e_1^*>0,
 \qquad q=\Abulk^{-1}e_1e^{i\theta}.
 \label{eq:anchored-response}
\end{equation}
Every hidden principal coordinate block of $\Abulk$ is positive and
invertible independently of any observation nonvanishing.
\end{proposition}
\begin{proof}
Zero quadratic energy would give $Bv=0$ and $v_1=0$, hence $v=0$
by connectivity. Also $\Abulk q=e_1q_1=e_1e^{i\theta}$.
A principal block has the energy of the corresponding vector extended
by zeros, so its positivity follows from that of the full matrix.
\end{proof}

For a coordinate split $v=(v_o,v_h)$, write the source equation as
\[
 \begin{pmatrix}A_{oo}&A_{oh}\\A_{ho}&A_{hh}\end{pmatrix}
 \begin{pmatrix}v_o\\v_h\end{pmatrix}
 =\begin{pmatrix}f_o\\f_h\end{pmatrix},\qquad
 f=e_1e^{i\theta}.
\]
The proved hidden-block inverse gives the exact reduction
\begin{align}
 v_h&=A_{hh}^{-1}(f_h-A_{ho}v_o),\nonumber\\
 (A_{oo}-A_{oh}A_{hh}^{-1}A_{ho})v_o
    &=f_o-A_{oh}A_{hh}^{-1}f_h.
 \label{eq:schur-source}
\end{align}
Thus both the effective operator and the induced source retain the
hidden-coordinate coupling. This is finite Schur elimination; for a
modern operator formulation see \cite{feshbach}. It concerns coordinate
blocks of the positive bulk operator, and supplies no lower bound for
the separate oscillatory observation $C$ below.
For elimination at coordinate $1$, the hidden block of
$B^*B$ is already positive by the same zero-extension argument. Its
unanchored source Schur complement is zero; adding the anchor changes
it to one. For example,
\begin{equation}
 B=\begin{pmatrix}-2&1&0\\0&-3&1\end{pmatrix},\quad
 q=(1,2,6)^T,\quad
 \Abulk=\begin{pmatrix}5&-2&0\\-2&10&-3\\0&-3&1\end{pmatrix}
 \label{eq:rational-source}
\end{equation}
has $\Abulk q=e_1$ and leading principal minors $5,46,1$.
The source construction assumes no zero-free observation. It differs
from a proposed elimination at a spectral value where the hidden inverse
has not been independently justified.

At fixed $t,\sigma$, put $U_Te_n=e^{iT\ell_n}e_n$.
Unitary multiplication of the edge row $(p,n)$ by $e^{iT\log(pn)}$
gives $B(T)U_T=U_EB(0)$, and therefore
\begin{equation}
 \Abulk(t,\sigma,T)=U_T\Abulk(t,\sigma,0)U_T^*.
 \label{eq:gauge}
\end{equation}
Its spectrum is independent of $T$ at fixed $t,\sigma$, while its
projected source response can cancel. Along physical $x$, $\sigma$
also changes. Positivity and a bulk gap thus supply no free observation
coercivity. Differentiation of the actual response does give
\begin{equation}
 q_x=i\theta_xq-\Abulk^{-1}(\partial_x\Abulk)q.
 \label{eq:source-derivative}
\end{equation}
The new proof object is this prescribed source response together with
its observation, rather than the positive spectrum alone.

\subsection{A real adjoint-transport certificate}

Realify $\mathbb C^N$ as $\mathbb R^{2N}$ with blocks
$(\Rea v_n,\Ima v_n)$, and denote by $B_R$ the realification of $B$.
Represent real covectors by column coefficient vectors. Define
\begin{align}
 Cv_R&=\left(\Rea\sum v_n,\;
    \frac1L\Rea\sum[-i\Omega+(-d+ic)\ell_n]v_n\right),\label{eq:C}\\
 b^Tv_R&=\Rea(e^{-i\theta}v_1),&
 a^Tv_R&=\Ima(e^{-i\theta}v_1).\nonumber
\end{align}
Then $Cq_R=(F/2,F'/(2L))$, $b^Tq_R=1$, $a^Tq_R=0$.
The latter anchor removes the free imaginary rotation of the complex
transport groundstate.

\begin{theorem}[Conditional anchored visibility certificate]
\label{thm:certificate}
Suppose real $y\in\mathbb R^2$, edge multipliers $\lambda$, a real scalar
$z_0$, and a residual covector $r$ satisfy
\begin{equation}
 b=C^Ty+B_R^T\lambda+z_0a+r.
 \label{eq:certificate}
\end{equation}
Write $r_n\in\mathbb R^2$ for its integer blocks and let
$\mathcal R=\sum_{n\le N}w_n\|r_n\|_2$.
If $\mathcal R<1$ and $y\ne0$, then
\begin{equation}
 \|Cq_R\|_2\ge\frac{1-\mathcal R}{\|y\|_2}.
 \label{eq:certificate-lower}
\end{equation}
Under Input~\ref{input:disk}, uniform outward cell bounds
$\mathcal R\le R_{\rm up}<1$, $\|y\|_2\le Y_{\rm up}$,
$\eta_N\le\eta_{\rm up}$ exclude every genuine joint zero whenever
\begin{equation}
 \frac{1-R_{\rm up}}{Y_{\rm up}}>\frac{\eta_{\rm up}}2.
 \label{eq:certificate-margin}
\end{equation}
Alternatively a sufficient directional margin is
\begin{equation}
 R_{\rm up}+\tfrac12\eta_{\rm up}
       \sup_{\rm cell}h(|y_2|,|y_1|)<1.
 \label{eq:directional-margin}
\end{equation}
\end{theorem}
\begin{proof}
Pair \eqref{eq:certificate} with $q_R$ and use both source anchors and
$B_Rq_R=0$. This gives $1=y\cdot Cq_R+r\cdot q_R$.
Since the $n$th state block has norm $w_n$,
$|r\cdot q_R|\le\mathcal R$, proving
\eqref{eq:certificate-lower} by Cauchy--Schwarz.
At a joint zero, Lemma~\ref{lem:paid-body} gives
$Cq_R=(\eta_N/2)(s,v)$ with $(s,v)\in\mathcal C_0$.
Because $v^2\le|v|$, its Euclidean norm is at most $\eta_N/2$.
This contradicts \eqref{eq:certificate-margin}; applying the exact
support function instead gives \eqref{eq:directional-margin}.
\end{proof}

The theorem proves a certificate lemma, not existence of effective
certificates. Unrestricted inversion that divides by the unknown pair
$(F,F')$ merely restates the nonvanishing question. A useful proof must
construct regular multipliers and bound their residuals without such
division. Using only the edges for $p=2,3$ leaves disconnected source
sectors; they must stay in the state and residual budget. The exact
transport constraints contain all phase channels and the complete cutoff;
a scalar block sign cannot replace them.

\begin{corollary}[Uniform shrinking-sector growth criterion]
\label{cor:growth}
Let $K\subset[1,3/2]$ be a closed interval. Suppose certificates
\eqref{eq:certificate} exist for the complete states throughout
\eqref{eq:sector} with $\kappa\in K$ and sufficiently small $t$, and
satisfy uniformly
\begin{equation}
 \mathcal R\le1-\delta,\qquad
 0<\|y\|_2\le C N^{\alpha_{\rm vis}},\qquad
 \delta>0,\quad
 \alpha_{\rm vis}<\inf_{\kappa\in K}\frac{\kappa+4}{8}.
 \label{eq:uniform-target}
\end{equation}
Then Input~\ref{input:disk} implies eventual joint-zero exclusion on
that subsector. The conservative sufficient exponent for all
$1\le\kappa\le3/2$ is $\alpha_{\rm vis}<5/8$.
\end{corollary}
\begin{proof}
Uniformly on the closed sector,
$\log N=\kappa/(2t)+o(1)$. Thus
\[
 \eta_N C N^{\alpha_{\rm vis}}
 \le5C\exp\left\{-\frac{\kappa(\kappa+4-8\alpha_{\rm vis})}{16t}
                         +o(1)\right\}\longrightarrow0.
\]
Eventually this is less than $2\delta$, establishing
\eqref{eq:certificate-margin}.
\end{proof}

The uniform hypotheses \eqref{eq:uniform-target} are open proof targets.
No multiplier theorem or asymptotic exclusion has been obtained.

\section{A continuous chord lift and its signed gradient source}
\label{sec:chord}

The continuous lift complements the doubled coherent arithmetic observations
and phase-space dynamics of program 13 \cite{microlocalnotes}.
For the genuine kernel let
\begin{equation}
 \mathfrak m_t(u)=e^{tu^2}\Phi_e(u),\qquad \psi_t(u)=\sqrt{\mathfrak m_t(u)},\qquad
 \partial_t\psi_t=\tfrac12u^2\psi_t.
 \label{eq:wave}
\end{equation}
Theta decay and its logarithmic derivative estimates give
$\psi_t\in\Sch(\mathbb R)$ with every fixed time derivative locally
uniformly on compact real time sets. Define the chord field and the
derivative-wavefunction overlap
\begin{align}
 \mathcal A_t(k,\ell)&=\int_{\mathbb R}e^{iku}
       \psi_t(u+\ell/2)\psi_t(u-\ell/2)\dd u,\label{eq:chord}\\
 \mathcal B_t(k,\ell)&=\int_{\mathbb R}e^{iku}
       \psi_t'(u+\ell/2)\psi_t'(u-\ell/2)\dd u.\nonumber
\end{align}
Both are real Schwartz fields on real $(k,\ell)$, and entire in $k$ for
each real $\ell$. No global complex-$\ell$ square-root branch is assumed.
The full fields retain pure-state coherence; the observed slice is
$\mathcal A_t(k,0)=2H_t(k)$.

\begin{proposition}[Chord and gradient-overlap equations]
\label{prop:chord}
Writing $\mathcal L=-\partial_k^2+\ell^2/4$ and
$\mathcal D_\perp=\partial_\ell^2+k^2/4$, one has
\begin{align}
 \partial_t\mathcal A&=\mathcal L\mathcal A,
 &\mathcal D_\perp\mathcal A&=-\mathcal B,\label{eq:chord-system}\\
 \partial_t\mathcal B&=\mathcal L\mathcal B
       -(k\partial_k+\ell\partial_\ell+1)\mathcal A.\nonumber
\end{align}
The transverse jets $C_{2r}=\partial_\ell^{2r}\mathcal A|_{\ell=0}$
satisfy
\begin{equation}
 \partial_tC_{2r}=-\partial_k^2C_{2r}
                         +\tfrac12r(2r-1)C_{2r-2}\quad(r\ge1).
 \label{eq:jet-hierarchy}
\end{equation}
\end{proposition}
\begin{proof}
The time multiplier of the product in \eqref{eq:chord} is
$(u+\ell/2)^2/2+(u-\ell/2)^2/2=u^2+\ell^2/4$.
Fourier transformation replaces $u^2$ by $-\partial_k^2$.
If $f=\psi_t(u+\ell/2)\psi_t(u-\ell/2)$, direct differentiation gives
$f_{\ell\ell}=f_{uu}/4-\psi_t'(u+\ell/2)\psi_t'(u-\ell/2)$.
Integration by parts proves the second identity. Finally
$[\mathcal D_\perp,\mathcal L]=k\partial_k+\ell\partial_\ell+1$;
apply $-\mathcal D_\perp$ to the first identity. Differentiating
$\ell^2\mathcal A/4$ $2r$ times at zero proves the hierarchy.
All integrations have vanishing Schwartz boundary terms.
\end{proof}

The observed slice is autonomous: transverse jets receive information
from it, but the hierarchy does not feed back into $H$. Any restriction
on its zeros must enter through the prepared theta state. Let
$C(k)=\mathcal A_{\ell\ell}(k,0)$ and use the Fourier convention
$\widehat f(k)=\int e^{iku}f(u)\dd u$. Then
\begin{equation}
 C=-k^2\mathcal A(k,0)/4-\widehat{|\psi_t'|^2}(k),
 \qquad |\psi_t'|^2=\mathfrak m_t\left(tu+\frac{\Phi_e'}{2\Phi_e}\right)^2.
 \label{eq:gradient-source}
\end{equation}
At a collision this identifies a signed full-theta gradient observation.
Positivity of its density gives no Fourier sign. With $W=-\log \mathfrak m_t$,
$C=-\widehat{\mathfrak m_tW''}/4$. If $C=P(\partial_k)\mathcal A(k,0)$ for a
finite constant-coefficient polynomial $P$, Fourier injectivity would
force $W''$ to be a polynomial in $u$. But the genuine tail gives
$W''(u)\sim16\pi e^{4|u|}$, excluding that particular finite local
closure. Nonlocal and arithmetic gluing relations remain possible.

\subsection{A bounded positive-threshold control}

The next estimate must distinguish genuine theta preparation from
generic purity and positive bulk norms.

\begin{proposition}[Pure-state lift with a positive double threshold]
\label{prop:control}
Put $t_*=1/40$ and
$\Phi_{\rm ctrl}(u)=(16u^4+24)e^{-(1+t_*)u^2}$.
For $0\le t\le1/20$, let $a_t=1+t_*-t>0$. Its even positive Schwartz
wavefunction satisfies \eqref{eq:chord-system} and has slice
\begin{equation}
 \mathcal A_t^{\rm ctrl}(k,0)=\sqrt\pi a_t^{-9/2}
 e^{-k^2/(4a_t)}
 [k^4-12a_tk^2+12a_t^2+24a_t^4].
 \label{eq:control}
\end{equation}
At $t=t_*$ the zeros are ordinary doubles at $k=\pm\sqrt6$.
The four polynomial zeros are nonreal for $t<t_*$ and all real for
$t\ge t_*$ in the stated interval.
\end{proposition}
\begin{proof}
Apply $16\partial_k^4+24$ to the Gaussian Fourier integral
$\sqrt{\pi/a_t}\,e^{-k^2/(4a_t)}$ to obtain
\eqref{eq:control}. At $a_t=1$ the polynomial is $(k^2-6)^2$.
Its discriminant in $k^2$ is $96a_t^2(1-a_t^2)$.
For $a_t>1$ its two roots in $k^2$ are nonreal. For $0<a_t\le1$ they
are $a_t(6\pm\sqrt{24(1-a_t^2)})$, both positive, establishing the
assertion. The wavefunction is the positive square root of its smooth
positive Gaussian density, so the preceding differential identities apply.
\end{proof}

This control has Gaussian tails and is a bounded-time heat family; it is
not the full all-time theta flow. It nevertheless rules out a claimed
local collision exclusion based only on the displayed chord equations,
rank-one purity, or positive phase-space norms. A useful new lemma must
use an actual theta coefficient or gluing property absent from this
control, and constrain the signed gradient overlap or the complete finite
current with all cross terms retained.

\section{Retained theta angular dynamics and gluing}
\label{sec:angular}

Before taking a rotor trace, retain the periodic angular field
\begin{equation}
 \Theta(r,y)=\sum_{n\in\mathbb Z}e^{-\pi n^2r+2\pi iny},
 \qquad k(u,y)=e^u[\Theta(e^{4u},y)-1],\quad u\ge0.
 \label{eq:angular}
\end{equation}
For $r>0$ its exact differential equation and Jacobi gluing are
\begin{align}
 \Theta_r&=\Theta_{yy}/(4\pi),
 &k_u&=k+e^{4u}k_{yy}/\pi,\label{eq:angular-pde}\\
 \Theta(r,y)&=r^{-1/2}\sum_{m\in\mathbb Z}
                        e^{-\pi(m-y)^2/r}.\label{eq:jacobi}
\end{align}
The first identities follow termwise; the last is Poisson/Jacobi
transformation \cite{dlmftheta}. The actual preparation has coefficient
one in every integer mode. Periodicity and the PDE alone permit other
initial Fourier coefficients, so they are not sufficient arithmetic
selection. Full coefficients and gluing must enter a new estimate.

For bounded complex parameters define
\begin{equation}
 \mathcal T_t(z,y)=\int_0^\infty e^{tu^2+izu}k(u,y)\dd u.
 \label{eq:Tfield}
\end{equation}
The nonzero theta modes have super-exponential $u$ decay, locally
uniformly with fixed derivatives in $y$; this justifies the integral,
its parameter derivatives, and $\partial_t\mathcal T=-\mathcal T_{zz}$.

\begin{proposition}[Angular differential-shift identity]
\label{prop:angular-shift}
The retained field obeys
\begin{equation}
 \frac1\pi\mathcal T_{yy}(z-4i,y)
 =-k(0,y)+2it\mathcal T_z(z,y)-(1+iz)\mathcal T(z,y).
 \label{eq:angular-shift}
\end{equation}
\end{proposition}
\begin{proof}
Multiply $k_u-k=e^{4u}k_{yy}/\pi$ by $e^{tu^2+izu}$ and integrate.
The left derivative integral is
$-k(0,y)-\int(2tu+iz)e^{tu^2+izu}k\dd u$.
Since $\mathcal T_z=i\int u e^{tu^2+izu}k\dd u$, subtraction of the
$k$ integral yields the right-hand side. The factor $e^{4u}$ is precisely
the shift $z\mapsto z-4i$. The boundary at infinity vanishes.
\end{proof}

At $y=0$, $k''-k=16\Phi$ and $k'(0)=-1$; differentiating the theta
series and using Jacobi evenness proves these identities. Hence
$J_t(x)=\Rea\mathcal T_t(x,0)$ is the usual auxiliary rotor cosine
readout, and two integrations by parts give
\begin{equation}
 16H_t(x)=1+(2t-1-x^2)J_t(x)+4txJ_t'(x)-4t^2J_t''(x).
 \label{eq:affine-ward}
\end{equation}
The endpoint constant is retained. The scalar endpoint identity itself
is shared by normalized positive Green-kernel controls and does not close
the sign \cite{thetanotes}; the retained angular field offers the fuller
arithmetic package to investigate.

The shift in \eqref{eq:angular-shift} is not covered by the imported
radius-$1/L$ disk at large height. Entireness does not pay for its use in
a finite approximation. For example, for real $y$ and $z=x-ib$,
absolute differentiation gives, for real $t,y,b$, a valid growth envelope
\begin{equation}
 |\partial_y^j\mathcal T_t(x-ib,y)|
 \le2(2\pi)^j\sum_{n\ge1}n^j
       \int_0^\infty e^{tu^2+(1+b)u-\pi n^2e^{4u}}\dd u.
 \label{eq:angular-growth}
\end{equation}
This is finite uniformly on compact $(t,b)$ sets. Its size is not a
small normalized approximation error. A source-to-observation estimate
using the shift must additionally control normalization, phase correlation,
and any finite-source remainder on that wider strip. This payment remains
an open analytical obligation.

\section{Logarithmic reflection and matched cutoff channels}
\label{sec:reflection}

For $P>0$ the map on $L^2(\mathbb R_+,\dd v)$
\begin{equation}
 (\mathcal I_Pf)(v)=\frac{\sqrt P}{v}f(P/v)
 \label{eq:inversion}
\end{equation}
is a linear unitary involution. Indeed substitution $w=P/v$ preserves
the norm, and applying the map twice gives $f$. The unitary logarithmic
coordinate change
\begin{equation}
 g(y)=P^{1/4}e^{y/2}f(\sqrt P\,e^y)
 \label{eq:logunitary}
\end{equation}
conjugates it to $g(y)\mapsto g(-y)$. Composition with complex
conjugation is a different, antiunitary map. A literal interval $[A,B]$
reflects to $[P/B,P/A]$.

For a carried chirp $f(v)=a(v)e^{iT\log v}$ with smooth local amplitude,
the Fourier phase $T\log v-2\pi kv$ has a positive-frequency saddle
$v_*=P/k$, where $P=T/(2\pi)>0$. Its leading local stationary coefficient
is
\begin{equation}
 e^{-i(T+\pi/4)}(\mathcal I_Pf)(k)
 =\frac{\sqrt P}{k}a(P/k)
          e^{i(T\log(P/k)-T-\pi/4)}.
 \label{eq:stationary-reflection}
\end{equation}
This identifies the coefficient by the usual stationary method
\cite{dlmfstationary}. It is not a uniform endpoint replacement or a
global remainder theorem. Such uses require exact finite Poisson integrals
or an endpoint-uniform approximation with an explicit global payment.

For the prescribed real weight $w(v)=e^{t\log^2v/4-\sigma\log v}$,
the modulus of \eqref{eq:logunitary} is proportional to $e^{ty^2/4}$
at the ideal center $\sigma=1/2+(t/4)\log P$.
It is even, but not globally square integrable at positive time; compact
blocks or justified weighted domains are required. The cutoff near
$N=\sqrt P$ is a reflection boundary at $y=0$.
In particular $[N/2,N]$ reflects approximately to $[N,2N]$, not to itself.
For a frozen physical point with $c\ne0$, put $\mu=\Omega/c$ and
$M_j=\sum_{n\le N}q_n(\log n-\mu)^j=X_j+iY_j$; let $\Gamma$ be a
frozen real coefficient. The ideal algebraic coefficient reflection
$M_j\mapsto(-1)^j\overline{M_j}$ fixes
$2Y_3^2+3X_2X_4-\Gamma X_2^2$ and supplies no opposite threshold sign
\cite{torusnotes,heat25}. Applying that ideal map to actual blocks still
requires payment for the center, carrier and weight mismatches.

The opportunity is a complete matched source or flux calculation:
resonant integrals, literal endpoints and half-weights, physical carrier
mismatch, and reflection defects together. These may generate regular
adjoint multipliers in \eqref{eq:certificate} or a theta-specific current
bound. The rational-frequency organization of the exact sublattice
integrals and the macroscopic local Fourier ceiling in Heat Note 24
\cite{heatnotes} provide an existing effective organization.
The lower-index complement, shifted
sublattices, both arithmetic channels and every mixed contribution must
remain in the selected representation. Alternate complete representations
are alternatives, not extra disjoint corrections to be added.

\section{A bounded proof program and its limits}
\label{sec:program}

The principal target is a theorem about the actual anchored source,
not a general positivity assertion on an enlarged space. Three bounded
steps are suitable for investigation.
\begin{enumerate}
 \item Construct a regular arithmetic adjoint-transport witness
 \eqref{eq:certificate} on one fixed-cutoff shrinking-sector cell, using
 literal prime transport and matched cutoff channels. Pay every residual
 block and multiplier, and aim at \eqref{eq:certificate-margin} or the
 directional margin. A cell excluded when the separately enclosed value
 and slope tests are inconclusive would demonstrate useful interference.
 A restricted-family failure should identify its precise residual or
 multiplier loss. Uniform estimates \eqref{eq:uniform-target}, if proved,
 would extend that mechanism to eventual subsector exclusion.
 \item Seek a full theta gradient/score correlation for
 \eqref{eq:chord-system}--\eqref{eq:gradient-source}. It must restrict a
 signed overlap or the complete current, and visibly distinguish the
 actual theta state from Proposition~\ref{prop:control}. Independent
 positive bulk energies, purity and absolute moment envelopes do not
 furnish that restriction.
 \item If the needed source correlation is absent after a scalar trace,
 use the integer coefficients and Jacobi gluing in
 \eqref{eq:angular-pde}--\eqref{eq:angular-shift}. First make the
 complex-shift growth and normalization payment effective. Another
 scalar endpoint Ward identity is not yet the required new theorem.
\end{enumerate}

The exact lifts provide a useful route to source visibility with a
first-jet budget, which avoids the fourth-jet payment of the ordinary-double
threshold route. They do not prove visibility. In particular, the positive
parent operator and its guaranteed hidden-block inverse are structural
facts, while the observed source response remains oscillatory. The chord
slice is autonomous, and its nonlocal theta constraints are still open.
Neither leading stationary reflection nor angular heat symmetry implies
uniform nonvanishing.

A successful certificate based on \eqref{eq:body} covers all multiplicities
on its stated domain. Global conclusions still require cutoff changes,
intervening cells, remaining shrinking sectors, complementary parameter
ranges, and the small-time endpoint. If one instead returns to an
ordinary-double fourth-jet test, higher deflation remains a separate
obligation. The preparation records and finite algebra controls in
\cite{heat25} verify identities and bookkeeping, not a huge-height sign
or the open multiplier hypotheses. No new RH conclusion follows.

\section*{Acknowledgements and preparation record}

This initial manuscript was drafted for Edward Baker with substantial
GPT-6.1-sol (Codex) LLM assistance, reasoning effort ultra, verified from
the recorded drafting-turn configuration. The displayed derivations,
source reading, algebra controls and parallel cross-audits are internal
LLM work, not independent mathematical validation. Classical heat-flow,
approximation, theta and stationary-phase inputs are credited below.
No mathematical novelty or priority claim is made. The standalone source
uses no external input, image or bibliography files.

\begin{thebibliography}{99}
\bibitem{debruijn}
N.~G.~de Bruijn, \emph{The roots of trigonometric integrals},
Duke Math. J. \textbf{17} (1950), 197--226.
\url{https://doi.org/10.1215/S0012-7094-50-01720-0}.

\bibitem{newman}
C.~M.~Newman, \emph{Fourier transforms with only real zeros},
Proc. Amer. Math. Soc. \textbf{61} (1976), 245--251.
\url{https://www.jstor.org/stable/2041319}.

\bibitem{polymath}
D.~H.~J.~Polymath, \emph{Effective approximation of heat flow evolution
of the Riemann $\xi$ function, and a new upper bound for the
de Bruijn--Newman constant}, arXiv:1904.12438v2 (2019).
\url{https://arxiv.org/abs/1904.12438}.

\bibitem{dlmftheta}
NIST Digital Library of Mathematical Functions, equation 20.7.32,
\emph{Theta modular transformations}.
\url{https://dlmf.nist.gov/20.7.E32}.

\bibitem{dlmfstationary}
NIST Digital Library of Mathematical Functions, \S2.3(iv),
\emph{Method of stationary phase}.
\url{https://dlmf.nist.gov/2.3#iv}.

\bibitem{feshbach}
G.~Dusson, I.~M.~Sigal, and B.~Stamm,
\emph{The Feshbach--Schur map and perturbation theory},
arXiv:2105.02058 (2021).
\url{https://arxiv.org/abs/2105.02058}.

\bibitem{heatnotes}
Edward Baker research project, with LLM assistance,
\emph{Newman collision reductions}, internal Heat Notes 8, 12, 14,
and 23--24, 9--10 October 2026. Repository source directory:
\texttt{papers/quasi-rh-exponent-descent/newman\_collisions/notes/}.
The complete disk interface and global threshold transfer remain imported
inputs from that investigation.

\bibitem{torusnotes}
Edward Baker research project, with LLM assistance,
\emph{Prime-phase torus signed reductions}, manuscript and internal
Notes 1--9, 10 October 2026. Repository source directory:
\texttt{newman\_collisions/09\_prime\_phase\_torus/}.

\bibitem{microlocalnotes}
Edward Baker research project, with LLM assistance,
\emph{Microlocal observations, coherent currents, and finite zero atlases
for the Riemann heat flow}, program 13 manuscript and internal notes,
10 October 2026. Repository source directory:
\texttt{newman\_collisions/13\_microlocal\_phase\_space/}.

\bibitem{thetanotes}
Edward Baker research project, with LLM assistance,
\emph{Theta lattice}, internal Notes 1--2, 10 October 2026.
Repository source directory:
\texttt{newman\_collisions/06\_theta\_lattice/}.

\bibitem{heat25}
Edward Baker research project, with LLM assistance,
\emph{Enlarged-state transport, chord dynamics, and theta gluing},
Heat Note 25, 10 October 2026, with its algebra replay and internal review.
Repository source:
\texttt{newman\_collisions/notes/25\_ENLARGED\_STATE\_TRANSPORT\_CHORD\_}
\texttt{DYNAMICS\_AND\_THETA\_GLUING\_20261010.md}.
\end{thebibliography}
\end{document}
